\documentclass[../main.tex]{subfiles}
\begin{document}
\chapter{Metric and Normed Spaces}
\section{Motivation}
Thus far in IA Analysis I and IB Analysis II, we have been doing analysis on $\R$.
To do this, we relied on the structure provided by $|\argdot|$ on $\R$ that allowed us to define convergence with unique limits.
We then looked at Cauchy sequences and we saw that all such sequences also converge due to the \textit{completeness} of $\R$.
We also saw that closed bounded intervals are ``special'' because of the Bolzano Weierstrass theorem.

We then moved onto functions $f: \R \to \R$, or more generally $f: E \subseteq \R \to \R$.
We then defined continuity and showed that closed bounded intervals are also ``special'' in this case because of the extreme value theorem.
Finally, in this course, we extended this to notions of uniform continuity.

These definitions and properties are quite useful and so we would like to extend these notions to functions $f: X \to \R$ or $f: X \to Y$.
But what such sets $X$ and $Y$ can we do this on, and how can we generalise the structures from $\R$?
\section{Metric Spaces}
\begin{remark}[Note]
  In this course, we will assume that all vector spaces are over $\R$.
\end{remark}
To impose some structure on a set $X$, we can introduce a \textit{metric}.
More specifically, when $X$ is a vector space, we do this through the notion of a \textit{norm}.
\begin{definition}[Metric]
  A \textit{metric} on a set $X$ is a ``distance'' function $d: X \times X \to \R$ satisfying the following for any $x, y, z \in X$: \label{metricDef}
  \begin{enumerate}
    \item $d(x, y) \geq 0$ (i.e. distance is non-negative) with $d(x, y) = 0 \iff x = y$ (i.e. two points are the same if and only if the distance between them is 0)
    \item $d(x, y) = d(y, x)$ (i.e. the distance from $x$ to $y$ is the same as from $y$ to $x$)
    \item $d(x, y) \leq d(x, z) + d(z, y)$ (i.e. the triangle inequality)
  \end{enumerate}
  We then call $(X, d)$ a \textit{metric space}.
\end{definition}
\begin{definition}[Norm]
  A \textit{norm} on a (real) vector space $V$ is a function $\norm{\argdot}: V \to \R$ satisfying the following for any $x, y \in V$ and $\lambda \in \R$: \label{normDef}
  \begin{enumerate}
    \item $\norm{x} \geq 0$ with equality if and only if $x = \vec{0}$.
    \item $\norm{\lambda x} = |\lambda| \norm{x}$
    \item $\norm{x + y} \leq \norm{x} + \norm{y}$
  \end{enumerate}
  We then call $(V, \norm{\argdot})$ a \textit{normed vector space} or \textit{normed space}.
\end{definition}
\begin{proposition}
  If $(V, \norm{\argdot})$ is a normed space and we define:
  \[
    d: V \times V \to \R,\ (x, y) \mapsto \norm{x - y}
  \]
  then $(V, d)$ is a metric space.
\end{proposition}
\begin{proof}
  Checking the conditions of \cref{metricDef}, we have:
  \begin{enumerate}
    \item By \cref{normDef} \textbf{(i)}, $d(x, y) = \norm{x - y} \geq 0$ with equality if and only if $x - y = \vec{0} \iff x = y$.
    \item By \cref{normDef} \textbf{(ii)}:
      \[
        d(x, y) = \norm{x - y} = \norm{-(y - x)} = |-1| \norm{y - x} = \norm{y - x} = d(x, y)
      \]
    \item By \cref{normDef} \textbf{(iii)}:
      \[
        d(x, y) = \norm{x - y} = \norm{x - z + z - y} \leq \norm{x - z} + \norm{z - y} = d(x, z) + d(z, y)
      \]
  \end{enumerate}
  So $d$ is indeed a metric on $V$.
\end{proof}
\begin{example}[Finite dimensional vector space]
  The prototypical example of a finite dimensional normed space is $\R^n$ equipped with various norms.

  Recall that:
  \[
    \R^n = \{(x_1, \ldots, x_n): x_j \in \R\ \forall j \in \{1, \ldots, n\}\}
  \]
  is a vector space under the usual component-wise addition and scalar multiplication.

  We can introduce the following norms:
  \begin{itemize}
    \item The \textit{Euclidean} or $\ell^2$-norm defined as:
      \[
        \norm{x}_{E} = \left[\sum_{i = 1}^{n} (x_i)^2\right]^{\frac{1}{2}}
      \]
      This is sometimes denoted $\norm{\argdot}_2$ or $\norm{\argdot}_{\ell^2}$.

      This is a norm on $\R^n$, \textbf{(i)} and \textbf{(ii)} are trivial to check, for \textbf{(iii)}, consider $x = (x_1, \ldots, x_n)$ and $y = (y_1, \ldots, y_n)$:
      \begin{align*}
        \norm{x + y}^2_E &= \sum_{i = 1}^{n} (x_i + y_i)^2 \\
                         &= \norm{x}^2_E + 2 x \cdot y+ \norm{y}^2_E \\
                         &\leq \norm{x}^2_E + 2 \norm{x}^2_E \norm{y}^2_E + \norm{y}^2_E \text{ by Cauchy Schwarz} \\
                         &= (\norm{x}^2_E + \norm{y}^2_E)^2 \implies \norm{x + y}_E \leq \norm{x}_E + \norm{y}_E
      \end{align*}
    \item The $\ell^1$-norm on $\R^n$ defined as:
      \[
        \norm{x}_{\ell^1} = \sum_{i = 1}^{n} |x_i|
      \]
      This is sometimes denoted $\norm{\argdot}_1$.
      Conditions \textbf{(i)} and \textbf{(ii)} are trivial and \textbf{(iii)} follows from the triangle inequality on $|\argdot|$.
    \item The $\ell^{\infty}$-norm on $\R^n$ defined as:
      \[
        \norm{x}_{\ell^{\infty}} = \sup \{|x_i| : i \in \{1, \ldots, n\}\}
      \]
      This is sometimes denoted $\norm{\argdot}_{\infty}$.
      Conditions \textbf{(i)} and \textbf{(ii)} follow from the properties of $|\argdot|$ on $\R$.
      For \textbf{(iii)}, we have:
      \begin{align*}
        \norm{x + y}_{\infty} &= \sup_{i \in \{1, \ldots, n\}} [|x_i + y_i|] \\
                              &\leq \sup_{i \in \{1, \ldots, n\}} [|x_i| + |y_i|] \text{ using the triangle inequality} \\
                              &\leq \sup_{i \in \{1, \ldots, n\}} |x_i| + \sup_{i \in \{1, \ldots, n\}} |y_i| \\
                              &= \norm{x}_{\infty} + \norm{y}_{\infty}
      \end{align*}
  \end{itemize}
\end{example}
\begin{remark}[Recap]
  Suppose that $V$ is an $\F$-vector space with subset $W \subseteq V$.
  Recall from IB Linear Algebra that the subspace tells us that if $\vec{0} \in W$ and for all $w_1, w_2 \in W$ and $\lambda \in \F$ we have $w_1 + \lambda w_2 \in W$, then $W$ is a subspace of $V$, denoted $W \leq V$.
\end{remark}
\begin{example}[Infinite dimensional vector space]
  Consider the vector space of real sequences, denoted $\R^{\N}$:
  \[
    \R^{\N} = \{\text{sequence $(x_k)_{k \in \N}$ with $x_k \in \R\ \forall k \in \N$}\}
  \]
  This forms a vector space with element-wise addition and scalar multiplication.
  In this case, we will define subspaces of $\R^\N$ that have a particular property and then use this property to define a norm on the subspace.
  \begin{itemize}
    \item Define the subset $\ell^1$ of $\R^{\N}$ as:
      \[
        \ell^1 = \left\{(x_k) \in \R^{\N}: \sum_{k = 1}^{\infty} |x_k| < \infty\right\}
      \]
      Clearly $\vec{0} \in \ell^1$, and since $|x_i + \lambda y_i| \leq |x_i| + |\lambda| |y_i|$, if $x, y \in \ell_1$ and $\lambda \in \R$, we have $x + \lambda y \in \ell^1$.
      Thus by the subspace test, $\ell^1 \leq \R^{\N}$.

      We then define the $\ell^1$-norm on the $\ell^1$ subspace as:
      \[
        \norm{x}_{1} = \sum_{i = 1}^{\infty} |x_i| \text{ for $x \in \ell^1$}
      \]
      The conditions for $\norm{\argdot}_1$ to be a norm are all inherited from $|\argdot|$ on $\R$.
    \item We can similarly define the subset $\ell^2$ as:
      \[
        \ell^2 = \left\{(x_k) \in \R^{\N}: \sum_{k = 1}^{\infty} x_k^2 < \infty\right\}
      \]
      Again, clearly $\vec{0} \in \ell^2$.
      For any $x, y \in \ell^2$ and $\lambda \in \R$, we see that:
      \[
        0 \leq (x_i - \lambda y_i)^2 = x^2_i + \lambda^2 y^2_i - 2\lambda x_i y_i \implies 2 \lambda x_i y_i \leq x^2_i + \lambda^2 y^2_i
      \]
      and so:
      \[
        (x_i + \lambda y_i)^2 = x^2_i + \lambda^2 y^2_i + 2\lambda x_i y_i \leq 2(x^2_i + \lambda^2 y^2_i)
      \]
      Thus $x + \lambda y \in \ell^2$ and so by the subspace test, we have $\ell^2 \leq \R^{\N}$.

      We then define the $\ell^2$-norm on the $\ell^2$ subspace as:
      \[
        \norm{x}_{2} = \left[\sum_{i = 1}^{\infty} x^2_i\right]^{\frac{1}{2}} \text{ for $x \in \ell^2$}
      \]
      Clearly conditions \textbf{(i)} and \textbf{(ii)} hold.
      For \textbf{(iii)}, we see that:
      \[
        \norm{x}_2 = \lim_{n \to \infty} \norm{(x_1, \ldots, x_n)}_{\ell^2}
      \]
      and so the triangle inequality for $\norm{\argdot}_2$ follows from the triangle inequality of $\norm{\argdot}_{\ell^2}$ on $\R^n$.
    \item We also define the subset $\ell^{\infty}$ as:
      \[
        \ell_{\infty} = \left\{(x_k) \in \R^{\N}: \sup_{k \in \N} |x_k| < \infty\right\}
      \]
      Again, clearly $\vec{0} \in \ell^{\infty}$.
      For any $x, y \in \ell^{\infty}$ and $\lambda \in \R$:
      \[
        \sup_{k \in \N} |x_k + \lambda y_k| \leq \sup_{k \in \N} |x_k| + |\lambda| \sup_{k \in \N} |y_k| < \infty
      \]
      so $x + \lambda y \in \ell^{\infty}$.
      Thus by the subspace test, we have $\ell^{\infty} \leq \R^{\N}$.

      We then define the $\ell^{\infty}$-norm on the $\ell^{\infty}$ space as:
      \[
        \norm{x}_{\infty} = \sup_{k \in \N} |x_k|
      \]
      Conditions \textbf{(i)} and \textbf{(ii)} are inherited from $|\argdot|$ on $\R$.
      For condition \textbf{(iii)}, we have:
      \[
        \norm{x}_{\infty} = \lim_{n \to \infty} \norm{(x_1, \ldots, x_n)}_{\ell^{\infty}}
      \]
      and so the triangle inequality for $\norm{\argdot}_{\infty}$ follows from the triangle inequality of $\norm{\argdot}_{\ell^{\infty}}$ on $\R^n$.
  \end{itemize}
\end{example}
\end{document}
